\chapter{The Corruption of Praxis}

\epigraph{``As the city declines from aristocracy to timocracy, from timocracy to
oligarchy, from oligarchy to democracy, from democracy to tyranny, so does the
soul.''}{\textcite{plato1997complete}, \emph{Republic} 543c--576b}

\section{The Cycle of Degeneration}

In the \emph{Republic}, Socrates traces the successive regimes by which a just
city degenerates into an unjust one. Each regime represents a different answer
to the question ``What is the ruling principle of the city?'' Honour
(timocracy), wealth (oligarchy), freedom (democracy), and tyranny are not
merely political forms but \emph{psychic} ones --- each corresponds to a
different structure of the soul, a different answer to the question ``What rules
the human being?''

We trace the corresponding degeneration of \praxis. When action forgets its
ground in \physis, it passes through four increasingly degraded regimes. Each is
a way of answering the question ``What governs action?'' when nature no longer
does.

\section{Timocracy of Action: Honour}

The first corruption is the regime of honour. Here, action is governed by the
desire for recognition. The practitioner acts not because nature requires it but
because others will see and approve. The goal is not the action's coherence but
its visibility.

The signs of the timocratic practitioner: they speak constantly of their goals
but rarely act on them. They optimise their public goal tree for impressiveness
rather than feasibility. They collect streaks not for the rhythm they sustain
but for the badge they display. They choose goals by their prestige --- ``I
will run a marathon'' --- rather than by their fit with the practitioner's
actual nature.

The timocratic regime is where the \praxis\ platform is most easily perverted.
The public-facing features --- streaks, achievements, partner ratings --- can
become ends in themselves, the puppets of the cave. The platform must be
designed to resist this, to make the practitioner's inward coherence more
visible to them than their outward display.

\section{Oligarchy of Action: Wealth}

The second corruption is the regime of wealth. Here, action is governed by the
desire for material gain. The practitioner acts only if the action produces a
measurable economic return. All other action is dismissed as inefficient or
irrelevant.

The signs of the oligarchic practitioner: they maintain only those goals that
have a clear financial payoff. They evaluate every action by its ROI. They
dismiss the non-economic domains of the goal tree --- spirit, community, play
--- as indulgences. They mistake the part of nature that produces wealth for
the whole of nature.

The oligarchic regime confronts the \praxis\ platform with its deepest
challenge. The platform must operate in a market economy; it must charge for
its services, and its users will naturally evaluate it in economic terms. The
platform's defence is the diversity of its goal domains: 14 life domains from
Body \& Fitness to Spirit \& Purpose, only two of which (Financial Security and
Wealth \& Assets) are economic. The very structure of the ontology asserts that
human action has more governors than money.

\section{Democracy of Action: Freedom}

The third corruption is the regime of freedom. Here, action is governed by
nothing at all --- or rather, by the proposition that all governors are equally
valid and none has authority. The practitioner acts on every impulse, pursues
every desire, declares every goal --- and completes none.

The signs of the democratic practitioner: their goal tree is a jungle of
competing commitments, none of which persists long enough to bear fruit. They
are seduced by every new possibility. They mistake the proliferation of goals
for the richness of life. Their coherence score oscillates wildly --- moments
of intense commitment followed by long periods of inertia.

The democratic regime is the most seductive because it presents itself as
freedom. The practitioner believes they are liberated from the constraints of
nature, free to choose any action at any time. But this is the freedom of the
shopping mall --- the freedom to choose among commodities, not the freedom to
live according to one's nature. The democratic practitioner is not free; they
are scattered, divided against themselves, the Fragmented Ego of the Mirror
Trap \parencite{macintyre1981virtue}.

\section{Tyranny of Action: Will}

The fourth corruption is the regime of tyranny. Here, action is governed by the
raw will of the practitioner --- the desire to impose their vision on the world
regardless of nature's constraints. The tyrannical practitioner does not ask
``what does nature require?'' but ``what do I will?'' And they will without
limit.

The signs of the tyrannical practitioner: they set goals that violate the
limits of their own body, their community, their ecosystem. They push past the
signals of failure as if failure were an enemy to be defeated rather than
information from nature. They treat the \physis\ engine's coherence score as
a score to be maximised by any means, including the manipulation of the score
itself. Their action becomes compulsive, driven, destructive --- both of
themselves and of the world around them.

The tyrannical regime is the endpoint of the degeneration. Here, \praxis\ has
entirely forgotten its ground in \physis. The tool has become the master. The
practitioner believes they are most powerful when they are most autonomous from
nature; in fact, they are most enslaved, for their will is not guided by any
intelligence beyond itself.

\section{The Mechanism of Decline}

How does \praxis\ decline through these regimes? Not by an external force but
by an internal logic. Each regime contains the seed of the next.

The practitioner who acts for honour (timocracy) discovers that honour is not a
durable good --- it depends on the opinions of others, which are fickle. So they
turn to something more solid: wealth (oligarchy). But wealth, once accumulated,
does not satisfy the need for meaning --- it only creates more desire. So they
turn to freedom (democracy), pursuing every desire indiscriminately. But this
scatters them, and the scattering creates a desperate need for unity. Into this
void steps the tyrant --- the will that unifies by dominating everything in its
path.

The mechanism is the same at each stage: \praxis\ seeks a governor, and each
false governor fails. Honour fails, wealth fails, freedom fails --- only the
tyrant's will seems to succeed, and even that success is the greatest failure
of all, for the tyrant is ruled by their own will, which has no limit and can
never be satisfied.

\section{The Antidote}

The antidote to this decline is not a fifth, better regime. It is the recognition
that the search for a human governor of action --- honour, wealth, freedom, or
will --- is itself the error. The only adequate governor of action is nature
itself (\physis), not because nature is stronger than the tyrant's will (it is),
but because nature is the source of the intelligibility of action. Only when
action is governed by nature does it have an end beyond itself, a criterion of
success that is not arbitrary, a ground that does not shift with every change
of fashion or fortune.

The \praxis\ platform and the \physis\ engine are built on this recognition.
They do not create the governance of action by nature; they reveal it. The
coherence score is not a human invention but a natural discovery --- the
translation of nature's own judgment into a form that practitioners can see and
learn from.

The decline into timocracy, oligarchy, democracy, and tyranny is not inevitable.
It is a consequence of forgetting that \praxis\ is the tool of \physis. The
practitioner who remembers this --- who cultivates the disciplines, attends to
coherence, and returns to witness for others --- has broken the cycle.

\begin{figure}[ht]
\centering
\begin{tikzpicture}[
  reg/.style={draw, rounded corners, align=center, minimum width=28mm,
              minimum height=8mm, font=\small},
  >={Stealth[length=2mm]}]
  \node[reg, fill=green!15] (praxis) {Praxis governed by Physis};

  \node[reg, fill=yellow!15, below=of praxis, yshift=-4mm] (timo)
    {Timocracy\\\tiny governed by Honour};
  \draw[->, thick] (praxis) -- node[right, font=\tiny] {forgetfulness} (timo);

  \node[reg, fill=orange!15, below=of timo, yshift=-4mm] (olig)
    {Oligarchy\\\tiny governed by Wealth};
  \draw[->, thick] (timo) -- node[right, font=\tiny] {honour fails} (olig);

  \node[reg, fill=red!15, below=of olig, yshift=-4mm] (demo)
    {Democracy\\\tiny governed by Freedom};
  \draw[->, thick] (olig) -- node[right, font=\tiny] {wealth fails} (demo);

  \node[reg, fill=red!30, below=of demo, yshift=-4mm] (tyr)
    {Tyranny\\\tiny governed by Will};
  \draw[->, thick] (demo) -- node[right, font=\tiny] {freedom scatters} (tyr);

  \draw[->, thick, cohgreen, dashed] (tyr.east) -- ++(20mm,0)
    node[right, font=\tiny, cohgreen] {Antidote: return to Physis};
\end{tikzpicture}
\caption{The fourfold corruption of \praxis. Each regime is governed by a false
principle that fails, opening the way to the next. The antidote is the return
to nature as the sole adequate governor.\label{fig:corruptions}}
\end{figure}
